\chapter{2015年山东卷（理）}

\section{单选题}

\begin{problem}
已知集合 \(A = \{x \mid x^2 - 4x + 3 < 0\}\)，\(B = \{x \mid 2 < x < 4\}\)，则 \(A \cap B = (\quad)\)

\choices
    {(1,3)}
    {(1,4)}
    {(2,3)}
    {(2,4)}
\end{problem}

\begin{problem}
若复数 \(z\) 满足 \(\frac{\bar z}{1 - \i} = \i\)，其中 \(\i\) 为虚数单位，则 \(z =\)（\quad）

\choices
    {\(1 - \i\)}
    {\(1 + \i\)}
    {\(-1 - \i\)}
    {\(-1 + \i\)}
\end{problem}

\begin{problem}
要得到函数 \( y = \sin \left(4x - \frac{\pi}{3}\right) \) 的图象，只需要将函数 \( y = \sin 4x \) 的图象（\quad）

\choices
    {向左平移 \(\frac{\pi}{12}\) 个单位}
    {向右平移 \(\frac{\pi}{12}\) 个单位}
    {向左平移 \(\frac{\pi}{3}\) 个单位}
    {向右平移 \(\frac{\pi}{3}\) 个单位}
\end{problem}

\begin{problem}
已知菱形 \(ABCD\) 的边长为 \(a\)，\(\angle ABC = 60^{\circ}\)，则 \(\overrightarrow{BD} \cdot \overrightarrow{CD} =\)（\quad）

\choices
    {\(-\frac{3}{2} a^{2}\)}
    {\(-\frac{3}{4} a^{2}\)}
    {\(\frac{3}{4} a^{2}\)}
    {\(\frac{3}{2} a^2\)}
\end{problem}

\begin{problem}
不等式 \(|x - 1| - |x - 5| < 2\) 的解集是（\quad）

\choices
    {\((- \infty, 4)\)}
    {\((- \infty, 1)\)}
    {(1, 4)}
    {(1, 5)}
\end{problem}

\begin{problem}
已知 \(x\)，\(y\) 满足约束条件 \(\begin{cases}
x - y \geqslant 0,\\
x + y \leqslant 2,\\
y \geqslant 0,
\end{cases}\) 若 \(z = ax + y\) 的最大值为 4，则 \(a =\)（\quad）

\choices
    {3}
    {2}
    {\(-2\)}
    {\(-3\)}
\end{problem}

\begin{problem}
在梯形 \(ABCD\) 中，\(\angle ABC = \frac{\pi}{2}\)，\(AD // BC\)，\(BC = 2AD = 2AB = 2\). 将梯形 \(ABCD\) 绕 \(AD\) 所在的直线旋转一周而形成的曲面所围成的几何体的体积为（\quad）

\choices
    {\(\frac{2\pi}{3}\)}
    {\( \frac{4\pi}{3} \)}
    {\( \frac{5\pi}{3} \)}
    {\(2\pi\)}
\end{problem}

\begin{problem}
已知某批零件的长度误差（单位：毫米）服从正态分布 \(N(0,3^2)\)，从中随机取一件，其长度误差落在区间 \((3,6)\) 内的概率为（\quad） （附：若随机变量 \( \xi \) 服从正态分布 \( N(\mu, \sigma^{2}) \)，则 \( P(\mu - \sigma < \xi < \mu + \sigma) = 68.26\% \)，\( P(\mu - 2\sigma < \xi < \mu + 2\sigma) = 95.44\% \) ）

\choices
    {\(4.56\%\)}
    {\(13.59\%\)}
    {\(27.18\%\)}
    {\(31.74\%\)}
\end{problem}

\begin{problem}
一条光线从点 \((-2, -3)\) 射出，经 \(y\) 轴反射后与圆 \((x + 3)^2 + (y - 2)^2 = 1\) 相切，则反射光线所在直线的斜率为（\quad）

\choices
    {\(-\frac{5}{3}\) 或 \(-\frac{3}{5}\)}
    {\(-\frac{3}{5}\) 或 \(-\frac{2}{3}\)}
    {\(-\frac{5}{4}\) 或 \(-\frac{4}{5}\)}
    {\(-\frac{4}{3}\) 或 \(-\frac{3}{4}\)}
\end{problem}

\begin{problem}
设函数 \( f(x) = \begin{cases}
3x - 1, & x < 1,\\
2^x, & x \geqslant 1,
\end{cases} \) 则满足 \( f[f(a)] = 2^{f(a)} \) 的 \( a \) 取值范围是（\quad）

\choices
    {\( \left[\frac{2}{3},1\right] \)}
    {\([0,1]\)}
    {\(\left[\frac{2}{3}, +\infty\right)\)}
    {\([1, +\infty)\)}
\end{problem}

\section{填空题}

\begin{problem}
观察下列各式： \(\text{C}_1^0 = 4^0\);\ \(\text{C}_3^0 + \text{C}_3^1 = 4^1\);\ \(\text{C}_5^0 + \text{C}_5^1 + \text{C}_5^2 = 4^2\);\ \(\text{C}_7^0 + \text{C}_7^1 + \text{C}_7^2 + \text{C}_7^3 = 4^3\);\ \(\cdots\). 照此规律，当 \(n \in \mathbf{N}^*\) 时，\(\text{C}_{2n-1}^0 + \text{C}_{2n-1}^1 + \text{C}_{2n-1}^2 + \cdots + \text{C}_{2n-1}^{n-1} = \)\fillinblank{}.
\end{problem}

\begin{problem}
若“\(\forall x \in \left[0, \frac{\pi}{4}\right]\)，\(\tan x \leqslant m\)”是真命题，则实数 \(m\) 的最小值为\fillinblank{}.
\end{problem}

\begin{problem}
执行如图所示的程序框图，输出的 \(T\) 的值为\fillinblank{}.

\texfigure[max width=0.42\linewidth]{img_repaint/2015/shandong_science/q13_fig1.tex}
\end{problem}

\begin{problem}
已知函数 \( f(x) = a^x + b \) (\(a > 0\)，\(a \neq 1\)) 的定义域和值域都是 \([-1, 0]\)，则 \( a + b = \)\fillinblank{}.
\end{problem}

\begin{problem}
平面直角坐标系 \(xOy\) 中，双曲线 \(C_1: \frac{x^2}{a^2} - \frac{y^2}{b^2} = 1 (a > 0, b > 0)\) 的渐近线与抛物线 \(C_2: x^2 = 2py (p > 0)\) 交于 \(O\)，\(A\)，\(B\). 若 \(\triangle OAB\) 的垂心为 \(C_2\) 的焦点，则 \(C_1\) 的离心率为\fillinblank{}.
\end{problem}

\section{解答题}

\begin{problem}
设 \( f(x) = \sin x \cos x - \cos^2 \left( x + \frac{\pi}{4} \right) \).
\begin{enumerate}
\item 求 \( f(x) \) 的单调区间；
\item 在锐角 \(\triangle ABC\) 中，角 \(A\)，\(B\)，\(C\) 的对边分别为 \(a\)，\(b\)，\(c\). 若 \(f\left(\frac{A}{2}\right) = 0\)，\(a = 1\)，求 \(\triangle ABC\) 面积的最大值.
\end{enumerate}
\end{problem}

\begin{problem}
如图，在三棱台 \(DEF - ABC\) 中，\(AB = 2DE\)，\(G\)，\(H\) 分别为 \(AC\)，\(BC\) 的中点.
\begin{enumerate}
\item 求证： \(BD //\) 平面 \(FGH\)；
\item 若 \(CF \perp\) 平面 \(ABC\)，\(AB \perp BC\)，\(CF = DE\)，\(\angle BAC = 45^\circ\)，求平面 \(FGH\) 与平面 \(ACFD\) 所成的角 （锐角） 的大小.
\end{enumerate}

\bitmapfigure[max width=0.42\linewidth]{img/2015/shandong_science/q17_fig1.png}
\end{problem}

\begin{problem}
设数列 \(\{a_{n}\}\) 的前 \(n\) 项和为 \(S_{n}\). 已知 \(2S_{n} = 3^{n} + 3\).
\begin{enumerate}
\item 求 \( \{a_{n}\} \) 的通项公式；
\item 若数列 \( \{b_{n}\} \) 满足 \( a_{n}b_{n}=\log_{3}a_{n} \)，求 \( \{b_{n}\} \) 的前 n 项和 \( T_{n} \).
\end{enumerate}
\end{problem}

\begin{problem}
若 \(n\) 是一个三位正整数，且 \(n\) 的个位数字大于十位数字，十位数字大于百位数字，则称 \(n\) 为“三位递增数”（如 137， 359， 567 等）. 在某次数学趣味活动中，每位参加者需从所有的“三位递增数”中随机抽取 1 个数，且只能抽取一次. 得分规则如下：若抽取的“三位递增数”的三个数字之积不能被 5 整除，参加者得 0 分；若能被 5 整除，但不能被 10 整除，得 \(-1\) 分；若能被 10 整除，得 1 分.
\begin{enumerate}
\item 写出所有个位数字是 5 的“三位递增数”；
\item 若甲参加活动，求甲得分 X 的分布列和数学期望 \(E(X)\).
\end{enumerate}
\end{problem}

\begin{problem}
平面直角坐标系 \(xOy\) 中，已知椭圆 \(C: \frac{x^2}{a^2} + \frac{y^2}{b^2} = 1 (a > b > 0)\) 的离心率为 \(\frac{\sqrt{3}}{2}\). 左，右焦点分别是 \(F_1\)，\(F_2\). 以 \(F_1\) 为圆心以 3 为半径的圆与以 \(F_2\) 为圆心以 1 为半径的圆相交，且交点在椭圆 \(C\) 上.
\begin{enumerate}
\item 求椭圆 \(C\) 的方程；
\item 设椭圆 \(E: \frac{x^2}{4a^2} + \frac{y^2}{4b^2} = 1\)，\(P\) 为椭圆 \(C\) 上任意一点，过点 \(P\) 的直线 \(y = kx + m\) 交椭圆 \(E\) 于 \(A\)，\(B\) 两点，射线 \(PO\) 交椭圆 \(E\) 于点 \(Q\). 
\begin{enumerate}
\item 求 \( \frac{|OQ|}{|OP|} \) 的值；
\item 求 \(\triangle ABQ\) 面积的最大值.
\end{enumerate}
\end{enumerate}
\end{problem}

\begin{problem}
设函数 \( f(x) = \ln (x + 1) + a(x^2 - x) \)，其中 \( a \in \R \).
\begin{enumerate}
\item 讨论函数 \( f(x) \) 极值点的个数，并说明理由；
\item 若 \(\forall x > 0\)，\(f(x) \geqslant 0\) 成立，求 \(a\) 的取值范围.
\end{enumerate}
\end{problem}
